% !TeX root = ../../Book2.tex
\section{多维振荡积分}
\label{b2:sec:2103}

多维情况下既可能完全没有驻点，也可能只在某些方向振荡。估计的有效维数取决于可利用的非退化方向数，而不单由积分变量的个数决定。

\subsection{相位函数}
\label{b2:subsec:210301}

若 \(\nabla\phi\ne0\)，定义
\[
 L=\frac{\nabla\phi\cdot\nabla}{2\pi\mathrm i\lambda|\nabla\phi|^2},
 \qquad Le^{2\pi\mathrm i\lambda\phi}=e^{2\pi\mathrm i\lambda\phi}.
\]
对紧支撑振幅反复转移 \(L\) 的形式伴随，可得任意阶 \(\lambda^{-N}\) 衰减。
展开伴随时出现的每一项都是相位有限阶导数、\(|\nabla\phi|^{-1}\) 的幂与振幅导数的乘积。在 \(|\nabla\phi|\ge c\) 的紧集上这些系数有界，故证明和一维完全相同，只是导数变成多重指标。

\ProseParagraphBreak
若积分域有边界，散度定理同时产生边界积分。比如平直立方体上的线性相位，即使内域完全无驻点，某些面上的切向梯度仍可能为零。这解释了为什么光滑截断与特征函数截断的 Fourier 衰减差别很大。

\subsection{振幅}
\label{b2:subsec:210302}

振幅支撑的大小、光滑性和消失阶都参与估计。给定 \(a_\lambda(x)=a(\lambda^\theta x)\)，其第 \(j\) 阶导数有 \(\lambda^{j\theta}\) 的增长；直接套用固定振幅定理会把这些增长藏进余项。先换变量 \(u=\lambda^\theta x\)，新的相位参数变为 \(\lambda^{1-2\theta}\)（二次相位），才看得清有效振荡尺度。

\ProseParagraphBreak
振幅可以包含较慢的参数，但必须统一控制其支撑和导数。一个常用情形是 \(a(x,z)\) 支撑于固定紧集，\(z\) 属于紧参数集，并且所需 \(x\) 导数一致有界。若相位也具有统一非退化下界，所有前述估计都可对 \(z\) 一致成立。参数集不紧时应另行核对这些界。

\subsection{曲率}
\label{b2:subsec:210303}

设光滑超曲面在局部写成 \(x_d=h(x')\)。曲面测度的 Fourier 变换在方向 \(\omega\) 上具有相位
\(\phi_\omega(x')=-\omega'\cdot x'-\omega_d h(x')\)。
驻点条件说明 \(\omega\) 平行于曲面的法向；驻点处 Hesse 矩阵 为
\(-\omega_dD^2h\)。在相应坐标中 Gaussian 曲率非零，等价于这个 Hesse 矩阵 可逆。因此紧支撑光滑曲面密度满足
\[
 |\widehat{a\,\mathrm d\sigma}(\xi)|\le C(1+|\xi|)^{-(d-1)/2}
\]
（假设支撑上 Gaussian 曲率处处非零）。
对方向单位球面作有限覆盖，分出驻点邻域与非驻点区域，驻相公式和非驻相法分别控制它们，便得到方向一致的常数。

\ProseParagraphBreak
若曲面是平面，沿法向相位是常数，便没有这样的衰减。若曲率只有 \(r\) 个非零方向，局部二次模型只有 \(r\) 个振荡变量，通常只能期待 \(|\xi|^{-r/2}\)。柱面是一个直接例子：圆周方向提供曲率，轴向仍可保持相干。仅说“曲面光滑”不足以推出球面型衰减。

\subsection{局部化}
\label{b2:subsec:210304}

局部化需要同时记录所选覆盖和每片的非退化下界。若 \(\chi_j\) 构成振幅支撑附近的有限光滑分割，则
\[
 I_\lambda(a,\phi)=\sum_j I_\lambda(a\chi_j,\phi).
\]
在非驻相片使用梯度下界，在驻相片使用 Hesse 矩阵 下界；余项的阶数可以统一选取，再将有限项相加。
对于一个逼近退化曲面的参数族，片数或下界可能变化。此时统一估计需要新的尺度分解，不能仅把每个固定参数的驻相公式相加。

\ProseParagraphBreak
以 \(\phi_s(t)=t^3-st\) 为例，\(s>0\) 时有两个非退化驻点，\(s=0\) 时合并为三阶驻点，\(s<0\) 时没有驻点。两个驻点的 Hesse 矩阵 在 \(s\downarrow0\) 时趋于零，分别使用二次公式的常数会发散。令 \(t=\lambda^{-1/3}u\)，可见转变发生在 \(s\lambda^{2/3}\) 为常数量级时。这个例子说明统一渐近式为何需要将参数与空间尺度一起组织。
